% Options for packages loaded elsewhere
\PassOptionsToPackage{unicode}{hyperref}
\PassOptionsToPackage{hyphens}{url}
%
\documentclass[
]{article}
\usepackage{amsmath,amssymb}
\usepackage{iftex}
\ifPDFTeX
  \usepackage[T1]{fontenc}
  \usepackage[utf8]{inputenc}
  \usepackage{textcomp} % provide euro and other symbols
\else % if luatex or xetex
  \usepackage{unicode-math} % this also loads fontspec
  \defaultfontfeatures{Scale=MatchLowercase}
  \defaultfontfeatures[\rmfamily]{Ligatures=TeX,Scale=1}
\fi
\usepackage{lmodern}
\ifPDFTeX\else
  % xetex/luatex font selection
\fi
% Use upquote if available, for straight quotes in verbatim environments
\IfFileExists{upquote.sty}{\usepackage{upquote}}{}
\IfFileExists{microtype.sty}{% use microtype if available
  \usepackage[]{microtype}
  \UseMicrotypeSet[protrusion]{basicmath} % disable protrusion for tt fonts
}{}
\makeatletter
\@ifundefined{KOMAClassName}{% if non-KOMA class
  \IfFileExists{parskip.sty}{%
    \usepackage{parskip}
  }{% else
    \setlength{\parindent}{0pt}
    \setlength{\parskip}{6pt plus 2pt minus 1pt}}
}{% if KOMA class
  \KOMAoptions{parskip=half}}
\makeatother
\usepackage{xcolor}
\usepackage[margin=1in]{geometry}
\usepackage{color}
\usepackage{fancyvrb}
\newcommand{\VerbBar}{|}
\newcommand{\VERB}{\Verb[commandchars=\\\{\}]}
\DefineVerbatimEnvironment{Highlighting}{Verbatim}{commandchars=\\\{\}}
% Add ',fontsize=\small' for more characters per line
\usepackage{framed}
\definecolor{shadecolor}{RGB}{248,248,248}
\newenvironment{Shaded}{\begin{snugshade}}{\end{snugshade}}
\newcommand{\AlertTok}[1]{\textcolor[rgb]{0.94,0.16,0.16}{#1}}
\newcommand{\AnnotationTok}[1]{\textcolor[rgb]{0.56,0.35,0.01}{\textbf{\textit{#1}}}}
\newcommand{\AttributeTok}[1]{\textcolor[rgb]{0.13,0.29,0.53}{#1}}
\newcommand{\BaseNTok}[1]{\textcolor[rgb]{0.00,0.00,0.81}{#1}}
\newcommand{\BuiltInTok}[1]{#1}
\newcommand{\CharTok}[1]{\textcolor[rgb]{0.31,0.60,0.02}{#1}}
\newcommand{\CommentTok}[1]{\textcolor[rgb]{0.56,0.35,0.01}{\textit{#1}}}
\newcommand{\CommentVarTok}[1]{\textcolor[rgb]{0.56,0.35,0.01}{\textbf{\textit{#1}}}}
\newcommand{\ConstantTok}[1]{\textcolor[rgb]{0.56,0.35,0.01}{#1}}
\newcommand{\ControlFlowTok}[1]{\textcolor[rgb]{0.13,0.29,0.53}{\textbf{#1}}}
\newcommand{\DataTypeTok}[1]{\textcolor[rgb]{0.13,0.29,0.53}{#1}}
\newcommand{\DecValTok}[1]{\textcolor[rgb]{0.00,0.00,0.81}{#1}}
\newcommand{\DocumentationTok}[1]{\textcolor[rgb]{0.56,0.35,0.01}{\textbf{\textit{#1}}}}
\newcommand{\ErrorTok}[1]{\textcolor[rgb]{0.64,0.00,0.00}{\textbf{#1}}}
\newcommand{\ExtensionTok}[1]{#1}
\newcommand{\FloatTok}[1]{\textcolor[rgb]{0.00,0.00,0.81}{#1}}
\newcommand{\FunctionTok}[1]{\textcolor[rgb]{0.13,0.29,0.53}{\textbf{#1}}}
\newcommand{\ImportTok}[1]{#1}
\newcommand{\InformationTok}[1]{\textcolor[rgb]{0.56,0.35,0.01}{\textbf{\textit{#1}}}}
\newcommand{\KeywordTok}[1]{\textcolor[rgb]{0.13,0.29,0.53}{\textbf{#1}}}
\newcommand{\NormalTok}[1]{#1}
\newcommand{\OperatorTok}[1]{\textcolor[rgb]{0.81,0.36,0.00}{\textbf{#1}}}
\newcommand{\OtherTok}[1]{\textcolor[rgb]{0.56,0.35,0.01}{#1}}
\newcommand{\PreprocessorTok}[1]{\textcolor[rgb]{0.56,0.35,0.01}{\textit{#1}}}
\newcommand{\RegionMarkerTok}[1]{#1}
\newcommand{\SpecialCharTok}[1]{\textcolor[rgb]{0.81,0.36,0.00}{\textbf{#1}}}
\newcommand{\SpecialStringTok}[1]{\textcolor[rgb]{0.31,0.60,0.02}{#1}}
\newcommand{\StringTok}[1]{\textcolor[rgb]{0.31,0.60,0.02}{#1}}
\newcommand{\VariableTok}[1]{\textcolor[rgb]{0.00,0.00,0.00}{#1}}
\newcommand{\VerbatimStringTok}[1]{\textcolor[rgb]{0.31,0.60,0.02}{#1}}
\newcommand{\WarningTok}[1]{\textcolor[rgb]{0.56,0.35,0.01}{\textbf{\textit{#1}}}}
\usepackage{longtable,booktabs,array}
\usepackage{calc} % for calculating minipage widths
% Correct order of tables after \paragraph or \subparagraph
\usepackage{etoolbox}
\makeatletter
\patchcmd\longtable{\par}{\if@noskipsec\mbox{}\fi\par}{}{}
\makeatother
% Allow footnotes in longtable head/foot
\IfFileExists{footnotehyper.sty}{\usepackage{footnotehyper}}{\usepackage{footnote}}
\makesavenoteenv{longtable}
\usepackage{graphicx}
\makeatletter
\def\maxwidth{\ifdim\Gin@nat@width>\linewidth\linewidth\else\Gin@nat@width\fi}
\def\maxheight{\ifdim\Gin@nat@height>\textheight\textheight\else\Gin@nat@height\fi}
\makeatother
% Scale images if necessary, so that they will not overflow the page
% margins by default, and it is still possible to overwrite the defaults
% using explicit options in \includegraphics[width, height, ...]{}
\setkeys{Gin}{width=\maxwidth,height=\maxheight,keepaspectratio}
% Set default figure placement to htbp
\makeatletter
\def\fps@figure{htbp}
\makeatother
\setlength{\emergencystretch}{3em} % prevent overfull lines
\providecommand{\tightlist}{%
  \setlength{\itemsep}{0pt}\setlength{\parskip}{0pt}}
\setcounter{secnumdepth}{-\maxdimen} % remove section numbering
\usepackage{booktabs}
\usepackage{longtable}
\usepackage{array}
\usepackage{multirow}
\usepackage{wrapfig}
\usepackage{float}
\usepackage{colortbl}
\usepackage{pdflscape}
\usepackage{tabu}
\usepackage{threeparttable}
\usepackage{threeparttablex}
\usepackage[normalem]{ulem}
\usepackage{makecell}
\usepackage{xcolor}
\ifLuaTeX
  \usepackage{selnolig}  % disable illegal ligatures
\fi
\usepackage{bookmark}
\IfFileExists{xurl.sty}{\usepackage{xurl}}{} % add URL line breaks if available
\urlstyle{same}
\hypersetup{
  pdftitle={project\_unit\_3},
  hidelinks,
  pdfcreator={LaTeX via pandoc}}

\title{project\_unit\_3}
\author{}
\date{\vspace{-2.5em}2025-04-16}

\begin{document}
\maketitle

Esta actividad tiene como fin desarrollar un modelo de regresión simple
con datos obtenidos del respositorio de datos UC Irvine. Para este
ejercicio se han utilizado las calificaciones de las asignatura de
Matemáticas y Portugués de dos escuelas de secundaria de Portugal. El
conjunto de datos no cuenta con datos perdidos.

En este informe se empieza por crear dos tablas que muestran las
variables cuantitativas que tiene cada conjunto de datos. Con esta
información se pueden empezar a establecer posibles correlaciones
presentes en las variables para posteriormente realizar el modelo de
regresión lineal.

Para este ejercicio solo se utilizarán las variables cualitativas para
realizar el modelo, ya que es una regresión lineal simple.

El problema a resolver en este ejercicio es la predicción de la nota
correspondiente al último o tercer periodo partiendo de las notas
anteriores, según el nivel de la correlación existente entre las
variables.

\section{Exploración de datos}\label{exploraciuxf3n-de-datos}

Este es el primer paso en cualquier proyecto, pues es necesario explorar
el conjunto de datos para conocer las variables que posee cada conjunto
de datos y el tipo de estas. En este caso, para practicidad, no se
manipula el merge hecho originalmente, sino que se usa cada una de las
bases para aplicar los modelos a cada base.

En este primer momento, se filtran las variables numéricas, ya que, como
se dijo al inicio, solo se usarán los datos cuantitativos en este
modelo.

\begin{Shaded}
\begin{Highlighting}[]
\NormalTok{student\_por\_num }\OtherTok{\textless{}{-}} \FunctionTok{select}\NormalTok{(student\_por, }\FunctionTok{where}\NormalTok{(is.numeric))}
\NormalTok{student\_mat\_num }\OtherTok{\textless{}{-}} \FunctionTok{select}\NormalTok{(student\_mat, }\FunctionTok{where}\NormalTok{(is.numeric))}

\FunctionTok{kable}\NormalTok{(}\FunctionTok{summary}\NormalTok{(student\_por\_num))}
\end{Highlighting}
\end{Shaded}

\begin{longtable}[]{@{}
  >{\raggedright\arraybackslash}p{(\columnwidth - 32\tabcolsep) * \real{0.0133}}
  >{\raggedright\arraybackslash}p{(\columnwidth - 32\tabcolsep) * \real{0.0619}}
  >{\raggedright\arraybackslash}p{(\columnwidth - 32\tabcolsep) * \real{0.0619}}
  >{\raggedright\arraybackslash}p{(\columnwidth - 32\tabcolsep) * \real{0.0619}}
  >{\raggedright\arraybackslash}p{(\columnwidth - 32\tabcolsep) * \real{0.0619}}
  >{\raggedright\arraybackslash}p{(\columnwidth - 32\tabcolsep) * \real{0.0619}}
  >{\raggedright\arraybackslash}p{(\columnwidth - 32\tabcolsep) * \real{0.0664}}
  >{\raggedright\arraybackslash}p{(\columnwidth - 32\tabcolsep) * \real{0.0619}}
  >{\raggedright\arraybackslash}p{(\columnwidth - 32\tabcolsep) * \real{0.0575}}
  >{\raggedright\arraybackslash}p{(\columnwidth - 32\tabcolsep) * \real{0.0619}}
  >{\raggedright\arraybackslash}p{(\columnwidth - 32\tabcolsep) * \real{0.0619}}
  >{\raggedright\arraybackslash}p{(\columnwidth - 32\tabcolsep) * \real{0.0575}}
  >{\raggedright\arraybackslash}p{(\columnwidth - 32\tabcolsep) * \real{0.0619}}
  >{\raggedright\arraybackslash}p{(\columnwidth - 32\tabcolsep) * \real{0.0664}}
  >{\raggedright\arraybackslash}p{(\columnwidth - 32\tabcolsep) * \real{0.0575}}
  >{\raggedright\arraybackslash}p{(\columnwidth - 32\tabcolsep) * \real{0.0619}}
  >{\raggedright\arraybackslash}p{(\columnwidth - 32\tabcolsep) * \real{0.0619}}@{}}
\toprule\noalign{}
\begin{minipage}[b]{\linewidth}\raggedright
\end{minipage} & \begin{minipage}[b]{\linewidth}\raggedright
age
\end{minipage} & \begin{minipage}[b]{\linewidth}\raggedright
Medu
\end{minipage} & \begin{minipage}[b]{\linewidth}\raggedright
Fedu
\end{minipage} & \begin{minipage}[b]{\linewidth}\raggedright
traveltime
\end{minipage} & \begin{minipage}[b]{\linewidth}\raggedright
studytime
\end{minipage} & \begin{minipage}[b]{\linewidth}\raggedright
failures
\end{minipage} & \begin{minipage}[b]{\linewidth}\raggedright
famrel
\end{minipage} & \begin{minipage}[b]{\linewidth}\raggedright
freetime
\end{minipage} & \begin{minipage}[b]{\linewidth}\raggedright
goout
\end{minipage} & \begin{minipage}[b]{\linewidth}\raggedright
Dalc
\end{minipage} & \begin{minipage}[b]{\linewidth}\raggedright
Walc
\end{minipage} & \begin{minipage}[b]{\linewidth}\raggedright
health
\end{minipage} & \begin{minipage}[b]{\linewidth}\raggedright
absences
\end{minipage} & \begin{minipage}[b]{\linewidth}\raggedright
G1
\end{minipage} & \begin{minipage}[b]{\linewidth}\raggedright
G2
\end{minipage} & \begin{minipage}[b]{\linewidth}\raggedright
G3
\end{minipage} \\
\midrule\noalign{}
\endhead
\bottomrule\noalign{}
\endlastfoot
& Min. :15.00 & Min. :0.000 & Min. :0.000 & Min. :1.000 & Min. :1.000 &
Min. :0.0000 & Min. :1.000 & Min. :1.00 & Min. :1.000 & Min. :1.000 &
Min. :1.00 & Min. :1.000 & Min. : 0.000 & Min. : 0.0 & Min. : 0.00 &
Min. : 0.00 \\
& 1st Qu.:16.00 & 1st Qu.:2.000 & 1st Qu.:1.000 & 1st Qu.:1.000 & 1st
Qu.:1.000 & 1st Qu.:0.0000 & 1st Qu.:4.000 & 1st Qu.:3.00 & 1st
Qu.:2.000 & 1st Qu.:1.000 & 1st Qu.:1.00 & 1st Qu.:2.000 & 1st Qu.:
0.000 & 1st Qu.:10.0 & 1st Qu.:10.00 & 1st Qu.:10.00 \\
& Median :17.00 & Median :2.000 & Median :2.000 & Median :1.000 & Median
:2.000 & Median :0.0000 & Median :4.000 & Median :3.00 & Median :3.000 &
Median :1.000 & Median :2.00 & Median :4.000 & Median : 2.000 & Median
:11.0 & Median :11.00 & Median :12.00 \\
& Mean :16.74 & Mean :2.515 & Mean :2.307 & Mean :1.569 & Mean :1.931 &
Mean :0.2219 & Mean :3.931 & Mean :3.18 & Mean :3.185 & Mean :1.502 &
Mean :2.28 & Mean :3.536 & Mean : 3.659 & Mean :11.4 & Mean :11.57 &
Mean :11.91 \\
& 3rd Qu.:18.00 & 3rd Qu.:4.000 & 3rd Qu.:3.000 & 3rd Qu.:2.000 & 3rd
Qu.:2.000 & 3rd Qu.:0.0000 & 3rd Qu.:5.000 & 3rd Qu.:4.00 & 3rd
Qu.:4.000 & 3rd Qu.:2.000 & 3rd Qu.:3.00 & 3rd Qu.:5.000 & 3rd Qu.:
6.000 & 3rd Qu.:13.0 & 3rd Qu.:13.00 & 3rd Qu.:14.00 \\
& Max. :22.00 & Max. :4.000 & Max. :4.000 & Max. :4.000 & Max. :4.000 &
Max. :3.0000 & Max. :5.000 & Max. :5.00 & Max. :5.000 & Max. :5.000 &
Max. :5.00 & Max. :5.000 & Max. :32.000 & Max. :19.0 & Max. :19.00 &
Max. :19.00 \\
\end{longtable}

\begin{Shaded}
\begin{Highlighting}[]
\FunctionTok{kable}\NormalTok{(}\FunctionTok{summary}\NormalTok{(student\_mat\_num))}
\end{Highlighting}
\end{Shaded}

\begin{longtable}[]{@{}
  >{\raggedright\arraybackslash}p{(\columnwidth - 32\tabcolsep) * \real{0.0132}}
  >{\raggedright\arraybackslash}p{(\columnwidth - 32\tabcolsep) * \real{0.0570}}
  >{\raggedright\arraybackslash}p{(\columnwidth - 32\tabcolsep) * \real{0.0614}}
  >{\raggedright\arraybackslash}p{(\columnwidth - 32\tabcolsep) * \real{0.0614}}
  >{\raggedright\arraybackslash}p{(\columnwidth - 32\tabcolsep) * \real{0.0614}}
  >{\raggedright\arraybackslash}p{(\columnwidth - 32\tabcolsep) * \real{0.0614}}
  >{\raggedright\arraybackslash}p{(\columnwidth - 32\tabcolsep) * \real{0.0658}}
  >{\raggedright\arraybackslash}p{(\columnwidth - 32\tabcolsep) * \real{0.0614}}
  >{\raggedright\arraybackslash}p{(\columnwidth - 32\tabcolsep) * \real{0.0614}}
  >{\raggedright\arraybackslash}p{(\columnwidth - 32\tabcolsep) * \real{0.0614}}
  >{\raggedright\arraybackslash}p{(\columnwidth - 32\tabcolsep) * \real{0.0614}}
  >{\raggedright\arraybackslash}p{(\columnwidth - 32\tabcolsep) * \real{0.0614}}
  >{\raggedright\arraybackslash}p{(\columnwidth - 32\tabcolsep) * \real{0.0614}}
  >{\raggedright\arraybackslash}p{(\columnwidth - 32\tabcolsep) * \real{0.0658}}
  >{\raggedright\arraybackslash}p{(\columnwidth - 32\tabcolsep) * \real{0.0614}}
  >{\raggedright\arraybackslash}p{(\columnwidth - 32\tabcolsep) * \real{0.0614}}
  >{\raggedright\arraybackslash}p{(\columnwidth - 32\tabcolsep) * \real{0.0614}}@{}}
\toprule\noalign{}
\begin{minipage}[b]{\linewidth}\raggedright
\end{minipage} & \begin{minipage}[b]{\linewidth}\raggedright
age
\end{minipage} & \begin{minipage}[b]{\linewidth}\raggedright
Medu
\end{minipage} & \begin{minipage}[b]{\linewidth}\raggedright
Fedu
\end{minipage} & \begin{minipage}[b]{\linewidth}\raggedright
traveltime
\end{minipage} & \begin{minipage}[b]{\linewidth}\raggedright
studytime
\end{minipage} & \begin{minipage}[b]{\linewidth}\raggedright
failures
\end{minipage} & \begin{minipage}[b]{\linewidth}\raggedright
famrel
\end{minipage} & \begin{minipage}[b]{\linewidth}\raggedright
freetime
\end{minipage} & \begin{minipage}[b]{\linewidth}\raggedright
goout
\end{minipage} & \begin{minipage}[b]{\linewidth}\raggedright
Dalc
\end{minipage} & \begin{minipage}[b]{\linewidth}\raggedright
Walc
\end{minipage} & \begin{minipage}[b]{\linewidth}\raggedright
health
\end{minipage} & \begin{minipage}[b]{\linewidth}\raggedright
absences
\end{minipage} & \begin{minipage}[b]{\linewidth}\raggedright
G1
\end{minipage} & \begin{minipage}[b]{\linewidth}\raggedright
G2
\end{minipage} & \begin{minipage}[b]{\linewidth}\raggedright
G3
\end{minipage} \\
\midrule\noalign{}
\endhead
\bottomrule\noalign{}
\endlastfoot
& Min. :15.0 & Min. :0.000 & Min. :0.000 & Min. :1.000 & Min. :1.000 &
Min. :0.0000 & Min. :1.000 & Min. :1.000 & Min. :1.000 & Min. :1.000 &
Min. :1.000 & Min. :1.000 & Min. : 0.000 & Min. : 3.00 & Min. : 0.00 &
Min. : 0.00 \\
& 1st Qu.:16.0 & 1st Qu.:2.000 & 1st Qu.:2.000 & 1st Qu.:1.000 & 1st
Qu.:1.000 & 1st Qu.:0.0000 & 1st Qu.:4.000 & 1st Qu.:3.000 & 1st
Qu.:2.000 & 1st Qu.:1.000 & 1st Qu.:1.000 & 1st Qu.:3.000 & 1st Qu.:
0.000 & 1st Qu.: 8.00 & 1st Qu.: 9.00 & 1st Qu.: 8.00 \\
& Median :17.0 & Median :3.000 & Median :2.000 & Median :1.000 & Median
:2.000 & Median :0.0000 & Median :4.000 & Median :3.000 & Median :3.000
& Median :1.000 & Median :2.000 & Median :4.000 & Median : 4.000 &
Median :11.00 & Median :11.00 & Median :11.00 \\
& Mean :16.7 & Mean :2.749 & Mean :2.522 & Mean :1.448 & Mean :2.035 &
Mean :0.3342 & Mean :3.944 & Mean :3.235 & Mean :3.109 & Mean :1.481 &
Mean :2.291 & Mean :3.554 & Mean : 5.709 & Mean :10.91 & Mean :10.71 &
Mean :10.42 \\
& 3rd Qu.:18.0 & 3rd Qu.:4.000 & 3rd Qu.:3.000 & 3rd Qu.:2.000 & 3rd
Qu.:2.000 & 3rd Qu.:0.0000 & 3rd Qu.:5.000 & 3rd Qu.:4.000 & 3rd
Qu.:4.000 & 3rd Qu.:2.000 & 3rd Qu.:3.000 & 3rd Qu.:5.000 & 3rd Qu.:
8.000 & 3rd Qu.:13.00 & 3rd Qu.:13.00 & 3rd Qu.:14.00 \\
& Max. :22.0 & Max. :4.000 & Max. :4.000 & Max. :4.000 & Max. :4.000 &
Max. :3.0000 & Max. :5.000 & Max. :5.000 & Max. :5.000 & Max. :5.000 &
Max. :5.000 & Max. :5.000 & Max. :75.000 & Max. :19.00 & Max. :19.00 &
Max. :20.00 \\
\end{longtable}

La tabla anterior muestra que existen 382 registros, los estudiantes
tienen entre 15 y 22 años, los tiempos de estudio varían entre 1 y 4
horas y que más alumnos tuvieron de nota cero en cada periodo académico,
comparado con la asginatura de portugués, como se ilustra a
continuación. De lo anterior se destaca que hubo estudiantes que
obtuvieron de calificación un cero (0). Estos datos, una vez se grafica,
son atípicos, por lo que se debe tomar una decisión sobre estos. Para
este caso, se considera que lo mejor es omitirlos para el análisis, no
eliminarlos. ESta decisión se sustenta en el hecho de que no son datos
inexistentes, sino que corresponden a calificaciones reales que pueden
ser porque el estudiante no presentó el examen o cualquier otra
consideración del o la docente del área.

Por lo anterior, ya que interesa trabajar con las calificaciones de los
estudiantes, se filtran las bases y se escogen solo las variables de
calificaciones y se eomiten los valores que estén por encima de los
límites superior e inferior de la base de la asignatura Portugués. Esto
último se realiza porque es el conjunto de datos que presenta estos
datos notablemente.

\begin{Shaded}
\begin{Highlighting}[]
\CommentTok{\# Limpieza}
\CommentTok{\# se filtran los datos de cero en matemáticas por se los outliers}
\NormalTok{student\_mat\_num\_filtrado }\OtherTok{\textless{}{-}} \FunctionTok{filter}\NormalTok{(student\_mat\_num, G1 }\SpecialCharTok{\textgreater{}} \DecValTok{0}\NormalTok{, G2 }\SpecialCharTok{\textgreater{}} \DecValTok{0}\NormalTok{, G3 }\SpecialCharTok{\textgreater{}} \DecValTok{0}\NormalTok{)}


\CommentTok{\# Calcular los límites}
\NormalTok{Q1 }\OtherTok{\textless{}{-}} \FunctionTok{quantile}\NormalTok{(student\_por\_num}\SpecialCharTok{$}\NormalTok{G3, }\FloatTok{0.25}\NormalTok{)}
\NormalTok{Q3 }\OtherTok{\textless{}{-}} \FunctionTok{quantile}\NormalTok{(student\_por\_num}\SpecialCharTok{$}\NormalTok{G3, }\FloatTok{0.75}\NormalTok{)}
\NormalTok{IQR\_value }\OtherTok{\textless{}{-}} \FunctionTok{IQR}\NormalTok{(student\_por\_num}\SpecialCharTok{$}\NormalTok{G3)}

\CommentTok{\# Definir límites para detectar atípicos}
\NormalTok{limite\_inferior }\OtherTok{\textless{}{-}}\NormalTok{ Q1 }\SpecialCharTok{{-}} \FloatTok{1.5} \SpecialCharTok{*}\NormalTok{ IQR\_value}
\NormalTok{limite\_superior }\OtherTok{\textless{}{-}}\NormalTok{ Q3 }\SpecialCharTok{+} \FloatTok{1.5} \SpecialCharTok{*}\NormalTok{ IQR\_value}

\CommentTok{\# Filtrar datos sin atípicos}
\NormalTok{student\_por\_num\_filtrado }\OtherTok{\textless{}{-}}\NormalTok{ student\_por\_num }\SpecialCharTok{\%\textgreater{}\%}
  \FunctionTok{filter}\NormalTok{(G3 }\SpecialCharTok{\textgreater{}=}\NormalTok{ limite\_inferior }\SpecialCharTok{\&}\NormalTok{ G3 }\SpecialCharTok{\textless{}=}\NormalTok{ limite\_superior)}
\end{Highlighting}
\end{Shaded}

En los pasos anteriores se han seleccionado solo los valores
cuantitativos, numéricos, y luego se filtraron en las calificaciones,
que son las variables de interés, para suprimir los valores atípicos en
este modelo. El siguiente paso es hacer una gráfica para entender el
coheficiente correlación de Pearson, distribución y linealidad que
tienen las variables entres sí. Con este paso, se puede observar si se
cumplen algunos requisitos, como el R cuadrado.

En síntesis, permite conocer si existe o no una relación entre las
variables, positiva o negativa, que lleva a establecer si cuando se
incrementa una variable se puede incrementar otra o disminuir; si el
comportamiento de de las variables es el mismo, que se incrementan o
disminuyen al mismo tiempo, o contrario, que una se incrementa y otra
disminuye.

\begin{Shaded}
\begin{Highlighting}[]
\CommentTok{\# Correlaciones}

\CommentTok{\# creación de conjuntos de cada set de datos con las calificaciones para ver correlaciones entre estas.}

\NormalTok{grades\_mat }\OtherTok{\textless{}{-}}\NormalTok{ student\_mat[student\_mat}\SpecialCharTok{$}\NormalTok{G1 }\SpecialCharTok{\textgreater{}} \DecValTok{0} \SpecialCharTok{\&}\NormalTok{ student\_mat}\SpecialCharTok{$}\NormalTok{G2 }\SpecialCharTok{\textgreater{}} \DecValTok{0} \SpecialCharTok{\&}\NormalTok{ student\_mat}\SpecialCharTok{$}\NormalTok{G3 }\SpecialCharTok{\textgreater{}} \DecValTok{0}\NormalTok{, }\FunctionTok{c}\NormalTok{(}\StringTok{"G1"}\NormalTok{, }\StringTok{"G2"}\NormalTok{, }\StringTok{"G3"}\NormalTok{)]}


\NormalTok{grades\_por }\OtherTok{\textless{}{-}}\NormalTok{ student\_por[student\_por}\SpecialCharTok{$}\NormalTok{G1 }\SpecialCharTok{\textgreater{}} \DecValTok{0} \SpecialCharTok{\&}\NormalTok{ student\_por}\SpecialCharTok{$}\NormalTok{G2 }\SpecialCharTok{\textgreater{}} \DecValTok{0} \SpecialCharTok{\&}\NormalTok{ student\_por}\SpecialCharTok{$}\NormalTok{G3 }\SpecialCharTok{\textgreater{}} \DecValTok{0}\NormalTok{, }\FunctionTok{c}\NormalTok{(}\StringTok{"G1"}\NormalTok{, }\StringTok{"G2"}\NormalTok{, }\StringTok{"G3"}\NormalTok{)]}


\CommentTok{\# Correlaciones}
\CommentTok{\# notas de matemáticas:}
\NormalTok{GGally}\SpecialCharTok{::}\FunctionTok{ggpairs}\NormalTok{(grades\_mat,}
                \AttributeTok{lower =} \FunctionTok{list}\NormalTok{(}\AttributeTok{continuous =} \StringTok{"smooth"}\NormalTok{),}
                \AttributeTok{diag =} \FunctionTok{list}\NormalTok{(}\AttributeTok{continuous =} \StringTok{"barDiag"}\NormalTok{),}
                \AttributeTok{upper =} \FunctionTok{list}\NormalTok{(}\AttributeTok{continuous =} \FunctionTok{wrap}\NormalTok{(}\StringTok{"cor"}\NormalTok{, }\AttributeTok{size =} \DecValTok{4}\NormalTok{)))}
\end{Highlighting}
\end{Shaded}

\begin{verbatim}
`stat_bin()` using `bins = 30`. Pick better value with `binwidth`.
`stat_bin()` using `bins = 30`. Pick better value with `binwidth`.
`stat_bin()` using `bins = 30`. Pick better value with `binwidth`.
\end{verbatim}

\includegraphics{proyecto_unidad_3_files/figure-latex/unnamed-chunk-3-1.pdf}

\begin{Shaded}
\begin{Highlighting}[]
\FunctionTok{boxplot}\NormalTok{(grades\_mat)}
\end{Highlighting}
\end{Shaded}

\includegraphics{proyecto_unidad_3_files/figure-latex/unnamed-chunk-3-2.pdf}

\begin{Shaded}
\begin{Highlighting}[]
\CommentTok{\# notas de portugués}
\NormalTok{GGally}\SpecialCharTok{::}\FunctionTok{ggpairs}\NormalTok{(grades\_por,}
                \AttributeTok{lower =} \FunctionTok{list}\NormalTok{(}\AttributeTok{continuous =} \StringTok{"smooth"}\NormalTok{),}
                \AttributeTok{diag =} \FunctionTok{list}\NormalTok{(}\AttributeTok{continuous =} \StringTok{"barDiag"}\NormalTok{),}
                \AttributeTok{upper =} \FunctionTok{list}\NormalTok{(}\AttributeTok{continuous =} \FunctionTok{wrap}\NormalTok{(}\StringTok{"cor"}\NormalTok{, }\AttributeTok{size =} \DecValTok{4}\NormalTok{)))}
\end{Highlighting}
\end{Shaded}

\begin{verbatim}
`stat_bin()` using `bins = 30`. Pick better value with `binwidth`.
`stat_bin()` using `bins = 30`. Pick better value with `binwidth`.
`stat_bin()` using `bins = 30`. Pick better value with `binwidth`.
\end{verbatim}

\includegraphics{proyecto_unidad_3_files/figure-latex/unnamed-chunk-3-3.pdf}

\begin{Shaded}
\begin{Highlighting}[]
\FunctionTok{boxplot}\NormalTok{(grades\_por)}
\end{Highlighting}
\end{Shaded}

\includegraphics{proyecto_unidad_3_files/figure-latex/unnamed-chunk-3-4.pdf}

Teniendo en cuenta lo anterior, se puede ver en que las variables G3 y
G2 tienen una correlación de 0,966 en la asignatura de Matemáticas; para
Portugués es de 0,934. En ambos conjuntos de datos hay una correlación
muy fuerte, de casi 1, lo que quiere decir que es casi exacta la
relación entre varibales (como lo exhiben los gráficos de puntos).

Otro aspecto a destacar de estos datos es que los histogranmas muestran
que hay una distribución que tiende a normal entre las variables,
relacionadas entre ellas. En cada variable adicional se puede evidenciar
un comportamiento similar (como se muestra en los anexos).

\section{Modelo de regresión lineal
simple}\label{modelo-de-regresiuxf3n-lineal-simple}

\begin{Shaded}
\begin{Highlighting}[]
\CommentTok{\# Modelo regresión lineal}
\CommentTok{\# Correlación de Pearson Matemáticas}

\FunctionTok{cor}\NormalTok{(student\_mat\_num\_filtrado}\SpecialCharTok{$}\NormalTok{G3, student\_mat\_num\_filtrado}\SpecialCharTok{$}\NormalTok{G1, }\AttributeTok{use =} \StringTok{"complete.obs"}\NormalTok{)}
\end{Highlighting}
\end{Shaded}

\begin{verbatim}
[1] 0.891805
\end{verbatim}

\begin{Shaded}
\begin{Highlighting}[]
\FunctionTok{cor}\NormalTok{(student\_mat\_num\_filtrado}\SpecialCharTok{$}\NormalTok{G3, student\_mat\_num\_filtrado}\SpecialCharTok{$}\NormalTok{G2, }\AttributeTok{use =} \StringTok{"complete.obs"}\NormalTok{)}
\end{Highlighting}
\end{Shaded}

\begin{verbatim}
[1] 0.9655825
\end{verbatim}

\begin{Shaded}
\begin{Highlighting}[]
\FunctionTok{cor}\NormalTok{(student\_mat\_num\_filtrado}\SpecialCharTok{$}\NormalTok{G3, student\_mat\_num\_filtrado}\SpecialCharTok{$}\NormalTok{studytime, }\AttributeTok{use =} \StringTok{"complete.obs"}\NormalTok{)}
\end{Highlighting}
\end{Shaded}

\begin{verbatim}
[1] 0.1267281
\end{verbatim}

\begin{Shaded}
\begin{Highlighting}[]
\CommentTok{\# Correlación de Pearson Portugués}
\FunctionTok{cor}\NormalTok{(student\_por\_num\_filtrado}\SpecialCharTok{$}\NormalTok{G3, student\_por\_num\_filtrado}\SpecialCharTok{$}\NormalTok{G1, }\AttributeTok{use =} \StringTok{"complete.obs"}\NormalTok{)}
\end{Highlighting}
\end{Shaded}

\begin{verbatim}
[1] 0.874777
\end{verbatim}

\begin{Shaded}
\begin{Highlighting}[]
\FunctionTok{cor}\NormalTok{(student\_por\_num\_filtrado}\SpecialCharTok{$}\NormalTok{G3, student\_por\_num\_filtrado}\SpecialCharTok{$}\NormalTok{G2, }\AttributeTok{use =} \StringTok{"complete.obs"}\NormalTok{)}
\end{Highlighting}
\end{Shaded}

\begin{verbatim}
[1] 0.9426912
\end{verbatim}

\begin{Shaded}
\begin{Highlighting}[]
\FunctionTok{cor}\NormalTok{(student\_por\_num\_filtrado}\SpecialCharTok{$}\NormalTok{G3, student\_por\_num\_filtrado}\SpecialCharTok{$}\NormalTok{studytime, }\AttributeTok{use =} \StringTok{"complete.obs"}\NormalTok{)}
\end{Highlighting}
\end{Shaded}

\begin{verbatim}
[1] 0.2498546
\end{verbatim}

\begin{Shaded}
\begin{Highlighting}[]
\CommentTok{\# Regresión matemáticas}
\NormalTok{modelo\_mat }\OtherTok{\textless{}{-}} \FunctionTok{lm}\NormalTok{(G3 }\SpecialCharTok{\textasciitilde{}}\NormalTok{ G2, }\AttributeTok{data =}\NormalTok{ student\_mat\_num\_filtrado)}
\FunctionTok{summary}\NormalTok{(modelo\_mat)}
\end{Highlighting}
\end{Shaded}

\begin{verbatim}

Call:
lm(formula = G3 ~ G2, data = student_mat_num_filtrado)

Residuals:
    Min      1Q  Median      3Q     Max 
-2.2172 -0.1978 -0.1494  0.8119  2.8022 

Coefficients:
            Estimate Std. Error t value Pr(>|t|)    
(Intercept)  0.27529    0.16686    1.65   0.0999 .  
G2           0.99031    0.01416   69.95   <2e-16 ***
---
Signif. codes:  0 '***' 0.001 '**' 0.01 '*' 0.05 '.' 0.1 ' ' 1

Residual standard error: 0.8407 on 355 degrees of freedom
Multiple R-squared:  0.9323,    Adjusted R-squared:  0.9322 
F-statistic:  4893 on 1 and 355 DF,  p-value: < 2.2e-16
\end{verbatim}

\begin{Shaded}
\begin{Highlighting}[]
\CommentTok{\#plot(modelo\_mat)}

\NormalTok{modelo\_por }\OtherTok{\textless{}{-}} \FunctionTok{lm}\NormalTok{(G3 }\SpecialCharTok{\textasciitilde{}}\NormalTok{ G2, }\AttributeTok{data =}\NormalTok{ student\_por\_num\_filtrado)}
\FunctionTok{summary}\NormalTok{(modelo\_por)}
\end{Highlighting}
\end{Shaded}

\begin{verbatim}

Call:
lm(formula = G3 ~ G2, data = student_por_num_filtrado)

Residuals:
    Min      1Q  Median      3Q     Max 
-3.6274 -0.4847 -0.2468  0.5629  5.4678 

Coefficients:
            Estimate Std. Error t value Pr(>|t|)    
(Intercept)  1.00793    0.16167   6.234  8.3e-10 ***
G2           0.95243    0.01342  70.969  < 2e-16 ***
---
Signif. codes:  0 '***' 0.001 '**' 0.01 '*' 0.05 '.' 0.1 ' ' 1

Residual standard error: 0.8872 on 631 degrees of freedom
Multiple R-squared:  0.8887,    Adjusted R-squared:  0.8885 
F-statistic:  5037 on 1 and 631 DF,  p-value: < 2.2e-16
\end{verbatim}

\begin{Shaded}
\begin{Highlighting}[]
\CommentTok{\#plot(modelo\_mat)}
\end{Highlighting}
\end{Shaded}

\section{Análisis de diagnóstico}\label{anuxe1lisis-de-diagnuxf3stico}

\begin{Shaded}
\begin{Highlighting}[]
\CommentTok{\# diagnóstico del modelo}

\FunctionTok{par}\NormalTok{(}\AttributeTok{mfrow =} \FunctionTok{c}\NormalTok{(}\DecValTok{2}\NormalTok{, }\DecValTok{2}\NormalTok{))}
\FunctionTok{plot}\NormalTok{(modelo\_mat, }\AttributeTok{col =} \StringTok{"blue"}\NormalTok{)}
\end{Highlighting}
\end{Shaded}

\includegraphics{proyecto_unidad_3_files/figure-latex/unnamed-chunk-5-1.pdf}

\begin{Shaded}
\begin{Highlighting}[]
\FunctionTok{par}\NormalTok{(}\AttributeTok{mfrow =} \FunctionTok{c}\NormalTok{(}\DecValTok{2}\NormalTok{, }\DecValTok{2}\NormalTok{))}
\FunctionTok{plot}\NormalTok{(modelo\_por, }\AttributeTok{col =} \StringTok{"red"}\NormalTok{)}
\end{Highlighting}
\end{Shaded}

\includegraphics{proyecto_unidad_3_files/figure-latex/unnamed-chunk-5-2.pdf}

\section{Análisis de datos
atípicos}\label{anuxe1lisis-de-datos-atuxedpicos}

\begin{Shaded}
\begin{Highlighting}[]
\CommentTok{\# Valores ajustados}
\NormalTok{ajuste\_mat }\OtherTok{\textless{}{-}}\NormalTok{ modelo\_mat}\SpecialCharTok{$}\NormalTok{fitted.values}

\NormalTok{ajuste\_por }\OtherTok{\textless{}{-}}\NormalTok{ modelo\_por}\SpecialCharTok{$}\NormalTok{fitted.values}

\CommentTok{\# Extracción de residuales}

\NormalTok{resid\_mat }\OtherTok{\textless{}{-}}\NormalTok{ modelo\_mat}\SpecialCharTok{$}\NormalTok{residuals}

\NormalTok{resid\_por }\OtherTok{\textless{}{-}}\NormalTok{ modelo\_por}\SpecialCharTok{$}\NormalTok{residuals}

\CommentTok{\# Predicciones}

\NormalTok{student\_mat\_num\_filtrado}\SpecialCharTok{$}\NormalTok{predicciones }\OtherTok{\textless{}{-}} \FunctionTok{predict}\NormalTok{(modelo\_mat)}

\NormalTok{student\_por\_num\_filtrado}\SpecialCharTok{$}\NormalTok{predicciones }\OtherTok{\textless{}{-}} \FunctionTok{predict}\NormalTok{(modelo\_por)}

\CommentTok{\# dataframe con los datos de cada modelo}
  \CommentTok{\#Matemáticas}
\NormalTok{resultados\_mat }\OtherTok{\textless{}{-}} \FunctionTok{data.frame}\NormalTok{(}
  \AttributeTok{G2 =}\NormalTok{ student\_mat\_num\_filtrado}\SpecialCharTok{$}\NormalTok{G2,}
  \AttributeTok{G3 =}\NormalTok{ student\_mat\_num\_filtrado}\SpecialCharTok{$}\NormalTok{G3,}
  \AttributeTok{predicciones =}\NormalTok{ ajuste\_mat,}
  \AttributeTok{residual =}\NormalTok{ resid\_mat}
\NormalTok{)}

\FunctionTok{print}\NormalTok{(}\FunctionTok{head}\NormalTok{(resultados\_mat))}
\end{Highlighting}
\end{Shaded}

\begin{verbatim}
  G2 G3 predicciones   residual
1  6  6     6.217169 -0.2171693
2  5  6     5.226855  0.7731447
3  8 10     8.197797  1.8022029
4 14 15    14.139681  0.8603192
5 10 10    10.178425 -0.1784250
6 15 15    15.129995 -0.1299948
\end{verbatim}

\begin{Shaded}
\begin{Highlighting}[]
  \CommentTok{\#Portugués}
\NormalTok{resultados\_por }\OtherTok{\textless{}{-}} \FunctionTok{data.frame}\NormalTok{(}
  \AttributeTok{G2 =}\NormalTok{ student\_por\_num\_filtrado}\SpecialCharTok{$}\NormalTok{G2,}
  \AttributeTok{G3 =}\NormalTok{ student\_por\_num\_filtrado}\SpecialCharTok{$}\NormalTok{G3,}
  \AttributeTok{predicciones =}\NormalTok{ ajuste\_por,}
  \AttributeTok{residual =}\NormalTok{ resid\_por}
\NormalTok{)}
\FunctionTok{print}\NormalTok{(}\FunctionTok{head}\NormalTok{(resultados\_por))}
\end{Highlighting}
\end{Shaded}

\begin{verbatim}
  G2 G3 predicciones   residual
1 11 11     11.48465 -0.4846542
2 11 11     11.48465 -0.4846542
3 13 12     13.38951 -1.3895139
4 14 14     14.34194 -0.3419438
5 13 13     13.38951 -0.3895139
6 12 13     12.43708  0.5629160
\end{verbatim}

\begin{Shaded}
\begin{Highlighting}[]
\CommentTok{\# recta de regresión}

  \CommentTok{\# Matemáticas}
\FunctionTok{ggplot}\NormalTok{(student\_mat\_num\_filtrado, }\FunctionTok{aes}\NormalTok{(}\AttributeTok{x =}\NormalTok{ G2, }\AttributeTok{y =}\NormalTok{ G3)) }\SpecialCharTok{+}
  \FunctionTok{geom\_smooth}\NormalTok{(}\AttributeTok{method =} \StringTok{"lm"}\NormalTok{, }\AttributeTok{se =} \ConstantTok{FALSE}\NormalTok{, }\AttributeTok{color =} \StringTok{"lightgreen"}\NormalTok{) }\SpecialCharTok{+} \CommentTok{\#color de la regresión}
  \FunctionTok{geom\_segment}\NormalTok{(}\FunctionTok{aes}\NormalTok{(}\AttributeTok{xend =}\NormalTok{ G1, }\AttributeTok{yend =}\NormalTok{ predicciones), }\AttributeTok{col =} \StringTok{"red"}\NormalTok{, }\AttributeTok{lty =} \StringTok{"dashed"}\NormalTok{) }\SpecialCharTok{+} \CommentTok{\# residuales línea roja}
  \FunctionTok{geom\_point}\NormalTok{() }\SpecialCharTok{+} \CommentTok{\#Puntos de datos observados}
  \FunctionTok{geom\_point}\NormalTok{(}\FunctionTok{aes}\NormalTok{(}\AttributeTok{y =}\NormalTok{ predicciones), }\AttributeTok{col =} \StringTok{"red"}\NormalTok{) }\SpecialCharTok{+} \CommentTok{\#puntos de predicción en rojo}
  \FunctionTok{theme\_light}\NormalTok{() }\CommentTok{\#mejora la visualización}
\end{Highlighting}
\end{Shaded}

\begin{verbatim}
`geom_smooth()` using formula = 'y ~ x'
\end{verbatim}

\includegraphics{proyecto_unidad_3_files/figure-latex/unnamed-chunk-6-1.pdf}

\begin{Shaded}
\begin{Highlighting}[]
  \CommentTok{\# Portugués}
\FunctionTok{ggplot}\NormalTok{(student\_por\_num\_filtrado, }\FunctionTok{aes}\NormalTok{(}\AttributeTok{x =}\NormalTok{ G2, }\AttributeTok{y =}\NormalTok{ G3)) }\SpecialCharTok{+}
  \FunctionTok{geom\_smooth}\NormalTok{(}\AttributeTok{method =} \StringTok{"lm"}\NormalTok{, }\AttributeTok{se =} \ConstantTok{FALSE}\NormalTok{, }\AttributeTok{color =} \StringTok{"lightblue"}\NormalTok{) }\SpecialCharTok{+} \CommentTok{\#color de la regresión}
  \FunctionTok{geom\_segment}\NormalTok{(}\FunctionTok{aes}\NormalTok{(}\AttributeTok{xend =}\NormalTok{ G1, }\AttributeTok{yend =}\NormalTok{ predicciones), }\AttributeTok{col =} \StringTok{"red"}\NormalTok{, }\AttributeTok{lty =} \StringTok{"dashed"}\NormalTok{) }\SpecialCharTok{+} \CommentTok{\# residuales línea roja}
  \FunctionTok{geom\_point}\NormalTok{() }\SpecialCharTok{+} \CommentTok{\#Puntos de datos observados}
  \FunctionTok{geom\_point}\NormalTok{(}\FunctionTok{aes}\NormalTok{(}\AttributeTok{y =}\NormalTok{ predicciones), }\AttributeTok{col =} \StringTok{"red"}\NormalTok{) }\SpecialCharTok{+} \CommentTok{\#puntos de predicción en rojo}
  \FunctionTok{theme\_light}\NormalTok{() }\CommentTok{\#mejora la visualización}
\end{Highlighting}
\end{Shaded}

\begin{verbatim}
`geom_smooth()` using formula = 'y ~ x'
\end{verbatim}

\includegraphics{proyecto_unidad_3_files/figure-latex/unnamed-chunk-6-2.pdf}

\begin{Shaded}
\begin{Highlighting}[]
\CommentTok{\# gráfico de residuos y valors ajustados}
  \CommentTok{\# Matemáticas}
\FunctionTok{ggplot}\NormalTok{(resultados\_mat, }\FunctionTok{aes}\NormalTok{(}\AttributeTok{x =}\NormalTok{ predicciones, }\AttributeTok{y =}\NormalTok{ residual)) }\SpecialCharTok{+}
  \FunctionTok{geom\_point}\NormalTok{() }\SpecialCharTok{+}
  \FunctionTok{geom\_hline}\NormalTok{(}\AttributeTok{yintercept =} \DecValTok{0}\NormalTok{, }\AttributeTok{linetype =} \StringTok{"dashed"}\NormalTok{, }\AttributeTok{color =} \StringTok{"red"}\NormalTok{) }\SpecialCharTok{+}
  \FunctionTok{labs}\NormalTok{(}\AttributeTok{title =} \StringTok{"Residuos vs. valores ajustados"}\NormalTok{, }\AttributeTok{x =} \StringTok{"Valores ajustados"}\NormalTok{, }\AttributeTok{y =} \StringTok{"Residuos"}\NormalTok{) }\SpecialCharTok{+}
  \FunctionTok{theme\_minimal}\NormalTok{()}
\end{Highlighting}
\end{Shaded}

\includegraphics{proyecto_unidad_3_files/figure-latex/unnamed-chunk-6-3.pdf}

\begin{Shaded}
\begin{Highlighting}[]
  \CommentTok{\# Portugués}
\FunctionTok{ggplot}\NormalTok{(resultados\_por, }\FunctionTok{aes}\NormalTok{(}\AttributeTok{x =}\NormalTok{ predicciones, }\AttributeTok{y =}\NormalTok{ residual)) }\SpecialCharTok{+}
  \FunctionTok{geom\_point}\NormalTok{() }\SpecialCharTok{+}
  \FunctionTok{geom\_hline}\NormalTok{(}\AttributeTok{yintercept =} \DecValTok{0}\NormalTok{, }\AttributeTok{linetype =} \StringTok{"dashed"}\NormalTok{, }\AttributeTok{color =} \StringTok{"red"}\NormalTok{) }\SpecialCharTok{+}
  \FunctionTok{labs}\NormalTok{(}\AttributeTok{title =} \StringTok{"Residuos vs. valores ajustados"}\NormalTok{, }\AttributeTok{x =} \StringTok{"Valores ajustados"}\NormalTok{, }\AttributeTok{y =} \StringTok{"Residuos"}\NormalTok{) }\SpecialCharTok{+}
  \FunctionTok{theme\_minimal}\NormalTok{()}
\end{Highlighting}
\end{Shaded}

\includegraphics{proyecto_unidad_3_files/figure-latex/unnamed-chunk-6-4.pdf}

\begin{Shaded}
\begin{Highlighting}[]
\CommentTok{\# Histograma de residuos}
  \CommentTok{\#Matemáticas}
\FunctionTok{qqnorm}\NormalTok{(resid\_mat)}
\FunctionTok{qqline}\NormalTok{(resid\_mat, }\AttributeTok{col =} \StringTok{"red"}\NormalTok{)}
\end{Highlighting}
\end{Shaded}

\includegraphics{proyecto_unidad_3_files/figure-latex/unnamed-chunk-6-5.pdf}

\begin{Shaded}
\begin{Highlighting}[]
  \CommentTok{\# Portugués}
\FunctionTok{qqnorm}\NormalTok{(resid\_por)}
\FunctionTok{qqline}\NormalTok{(resid\_por, }\AttributeTok{col =} \StringTok{"blue"}\NormalTok{)}
\end{Highlighting}
\end{Shaded}

\includegraphics{proyecto_unidad_3_files/figure-latex/unnamed-chunk-6-6.pdf}

\begin{Shaded}
\begin{Highlighting}[]
\CommentTok{\#test de normalidad}
  \CommentTok{\# Matemáticas}
\FunctionTok{shapiro.test}\NormalTok{(resid\_mat)}
\end{Highlighting}
\end{Shaded}

\begin{verbatim}

    Shapiro-Wilk normality test

data:  resid_mat
W = 0.89793, p-value = 9.682e-15
\end{verbatim}

\begin{Shaded}
\begin{Highlighting}[]
  \CommentTok{\# Portugués}
\FunctionTok{shapiro.test}\NormalTok{(resid\_por)}
\end{Highlighting}
\end{Shaded}

\begin{verbatim}

    Shapiro-Wilk normality test

data:  resid_por
W = 0.92957, p-value < 2.2e-16
\end{verbatim}

\begin{Shaded}
\begin{Highlighting}[]
\CommentTok{\#test de homocedasticidad}

\CommentTok{\# Matemáticas}
\FunctionTok{bptest}\NormalTok{(modelo\_mat)}
\end{Highlighting}
\end{Shaded}

\begin{verbatim}

    studentized Breusch-Pagan test

data:  modelo_mat
BP = 12.297, df = 1, p-value = 0.0004536
\end{verbatim}

\begin{Shaded}
\begin{Highlighting}[]
  \CommentTok{\# Portugués}
\FunctionTok{bptest}\NormalTok{(modelo\_por)}
\end{Highlighting}
\end{Shaded}

\begin{verbatim}

    studentized Breusch-Pagan test

data:  modelo_por
BP = 4.6101, df = 1, p-value = 0.03178
\end{verbatim}

\section{Intervalos de confianza}\label{intervalos-de-confianza}

\section{Discusión y conclusiones}\label{discusiuxf3n-y-conclusiones}

Aunque existe una correlación fuerte entre la variable G3 y las
variables G2 y G1, no se puede establecer el principio de normalidad de
Shapiro ni bptest, puesto que el número de variables que se han
utilizado en el ejemplo no permiten que se puedan presentar. Aunque en
el código que se usó previo a este documento se intentó hacer con
logaritmo y otros métodos, no se pudo cumplir este principio.

Ahora bien, por tratarse de un ejemplo académico y práctico, se deja tal
como está.

La correlación vista entre las notas se presenta porque corresponden a
las calificaciones de las asignaturas de cada periodo académico en las
escuelas analizadas (G1 es del primero, G2, del segundo, y G3, del
tercero). Esto ocurre, ya que no se podrían tener las notas del tercer
periodo si no se tienen las del primero ni las del segundo (Cortes y
Silva, 2008).

\section{Referencias
bibliográficas}\label{referencias-bibliogruxe1ficas}

Cortez, P. y Silva, A. (2008). Student Performance {[}Dataset{]}. UCI
Machine Learning Repository. \url{doi://doi.org/10.24432/C5TG7T}.

\section{Anexos}\label{anexos}

En este apartado se escribe todo el código que no se usó en el cuerpo
del informe por dos razones: la primera, los requisitos de extensión
para esta actividad no permiten que estén ahí; la segunda, no son tan
relevantes para el informe, aunque dan cuenta de algunas características
del conjunto de datos escogido para esta actividad.

El primer código que se escribe es aquel que permite observar las
variables cuantitativas escogidas del conjunto de datos que corresponde
a la asignatura de Portugués. Para esta variable se ha escogido usar el
color ``lightgreen'' para diferenciarla de Matemáticas. Las variables
escogidas son edad (age), tiempo dedicado a estudiar (studytime), el
corte 1 (G1), el corte 2 (G2) y le corte 3 (G3)

\begin{Shaded}
\begin{Highlighting}[]
\CommentTok{\# Portugués}
\NormalTok{por1 }\OtherTok{\textless{}{-}} \FunctionTok{ggplot}\NormalTok{(student\_por\_num\_filtrado, }\FunctionTok{aes}\NormalTok{(}\AttributeTok{x =}\NormalTok{ age)) }\SpecialCharTok{+} \FunctionTok{geom\_bar}\NormalTok{(}\AttributeTok{fill =} \StringTok{"lightgreen"}\NormalTok{, }\AttributeTok{color =} \StringTok{"black"}\NormalTok{) }\SpecialCharTok{+} \FunctionTok{theme\_minimal}\NormalTok{()}
\NormalTok{por2 }\OtherTok{\textless{}{-}} \FunctionTok{ggplot}\NormalTok{(student\_por\_num\_filtrado, }\FunctionTok{aes}\NormalTok{(}\AttributeTok{x =}\NormalTok{ studytime)) }\SpecialCharTok{+} \FunctionTok{geom\_bar}\NormalTok{(}\AttributeTok{fill =} \StringTok{"lightgreen"}\NormalTok{, }\AttributeTok{color =} \StringTok{"black"}\NormalTok{) }\SpecialCharTok{+} \FunctionTok{theme\_minimal}\NormalTok{()}
\NormalTok{por3 }\OtherTok{\textless{}{-}} \FunctionTok{ggplot}\NormalTok{(student\_por\_num\_filtrado, }\FunctionTok{aes}\NormalTok{(}\AttributeTok{x =}\NormalTok{ G1)) }\SpecialCharTok{+} \FunctionTok{geom\_bar}\NormalTok{(}\AttributeTok{fill =} \StringTok{"lightgreen"}\NormalTok{, }\AttributeTok{color =} \StringTok{"black"}\NormalTok{) }\SpecialCharTok{+} \FunctionTok{theme\_minimal}\NormalTok{()}
\NormalTok{por4 }\OtherTok{\textless{}{-}} \FunctionTok{ggplot}\NormalTok{(student\_por\_num\_filtrado, }\FunctionTok{aes}\NormalTok{(}\AttributeTok{x =}\NormalTok{ G2)) }\SpecialCharTok{+} \FunctionTok{geom\_bar}\NormalTok{(}\AttributeTok{fill =} \StringTok{"lightgreen"}\NormalTok{, }\AttributeTok{color =} \StringTok{"black"}\NormalTok{) }\SpecialCharTok{+} \FunctionTok{theme\_minimal}\NormalTok{()}
\NormalTok{por5 }\OtherTok{\textless{}{-}} \FunctionTok{ggplot}\NormalTok{(student\_por\_num\_filtrado, }\FunctionTok{aes}\NormalTok{(}\AttributeTok{x =}\NormalTok{ G3)) }\SpecialCharTok{+} \FunctionTok{geom\_bar}\NormalTok{(}\AttributeTok{fill =} \StringTok{"lightgreen"}\NormalTok{, }\AttributeTok{color =} \StringTok{"black"}\NormalTok{) }\SpecialCharTok{+} \FunctionTok{theme\_minimal}\NormalTok{()}

\NormalTok{(por1 }\SpecialCharTok{|}\NormalTok{ por2 }\SpecialCharTok{|}\NormalTok{ por3)}\SpecialCharTok{/}\NormalTok{(por4 }\SpecialCharTok{|}\NormalTok{ por5)}
\end{Highlighting}
\end{Shaded}

\includegraphics{proyecto_unidad_3_files/figure-latex/unnamed-chunk-8-1.pdf}

El siguiente código es el que permite, al igual que el anterior, conocer
las variables cuantitativas escogidas del conjunto de datos utilizado en
el presente ejercicio. Para este caso, ha sido de la asignatura de
Matemáticas, para la cual se usa el color ``lightblue''. Las variables
escogidas son edad (age), tiempo dedicado a estudiar (studytime), el
corte 1 (G1), el corte 2 (G2) y le corte 3 (G3)

\begin{Shaded}
\begin{Highlighting}[]
\CommentTok{\# Matemáticas}
\NormalTok{mat1 }\OtherTok{\textless{}{-}} \FunctionTok{ggplot}\NormalTok{(student\_mat\_num\_filtrado, }\FunctionTok{aes}\NormalTok{(}\AttributeTok{x =}\NormalTok{ age)) }\SpecialCharTok{+} \FunctionTok{geom\_bar}\NormalTok{(}\AttributeTok{fill =} \StringTok{"lightblue"}\NormalTok{, }\AttributeTok{color =} \StringTok{"black"}\NormalTok{) }\SpecialCharTok{+} \FunctionTok{theme\_minimal}\NormalTok{()}
\NormalTok{mat2 }\OtherTok{\textless{}{-}} \FunctionTok{ggplot}\NormalTok{(student\_mat\_num\_filtrado, }\FunctionTok{aes}\NormalTok{(}\AttributeTok{x =}\NormalTok{ studytime)) }\SpecialCharTok{+} \FunctionTok{geom\_bar}\NormalTok{(}\AttributeTok{fill =} \StringTok{"lightblue"}\NormalTok{, }\AttributeTok{color =} \StringTok{"black"}\NormalTok{) }\SpecialCharTok{+} \FunctionTok{theme\_minimal}\NormalTok{()}
\NormalTok{mat3 }\OtherTok{\textless{}{-}} \FunctionTok{ggplot}\NormalTok{(student\_mat\_num\_filtrado, }\FunctionTok{aes}\NormalTok{(}\AttributeTok{x =}\NormalTok{ G1)) }\SpecialCharTok{+} \FunctionTok{geom\_bar}\NormalTok{(}\AttributeTok{fill =} \StringTok{"lightblue"}\NormalTok{, }\AttributeTok{color =} \StringTok{"black"}\NormalTok{) }\SpecialCharTok{+} \FunctionTok{theme\_minimal}\NormalTok{()}
\NormalTok{mat4 }\OtherTok{\textless{}{-}} \FunctionTok{ggplot}\NormalTok{(student\_mat\_num\_filtrado, }\FunctionTok{aes}\NormalTok{(}\AttributeTok{x =}\NormalTok{ G2)) }\SpecialCharTok{+} \FunctionTok{geom\_bar}\NormalTok{(}\AttributeTok{fill =} \StringTok{"lightblue"}\NormalTok{, }\AttributeTok{color =} \StringTok{"black"}\NormalTok{) }\SpecialCharTok{+} \FunctionTok{theme\_minimal}\NormalTok{()}
\NormalTok{mat5 }\OtherTok{\textless{}{-}} \FunctionTok{ggplot}\NormalTok{(student\_mat\_num\_filtrado, }\FunctionTok{aes}\NormalTok{(}\AttributeTok{x =}\NormalTok{ G3)) }\SpecialCharTok{+} \FunctionTok{geom\_bar}\NormalTok{(}\AttributeTok{fill =} \StringTok{"lightblue"}\NormalTok{, }\AttributeTok{color =} \StringTok{"black"}\NormalTok{) }\SpecialCharTok{+} \FunctionTok{theme\_minimal}\NormalTok{()}

\NormalTok{(mat1 }\SpecialCharTok{|}\NormalTok{ mat2 }\SpecialCharTok{|}\NormalTok{ mat3)}\SpecialCharTok{/}\NormalTok{(mat4 }\SpecialCharTok{|}\NormalTok{ mat5)}
\end{Highlighting}
\end{Shaded}

\includegraphics{proyecto_unidad_3_files/figure-latex/unnamed-chunk-9-1.pdf}

Actividad:

Caso de estudio y creación de herramienta interactiva A pesar de los
avances en el nivel educativo de la población portuguesa en las últimas
décadas, las estadísticas continúan ubicando a Portugal entre los países
con peores indicadores en Europa, principalmente debido a sus elevadas
tasas de fracaso escolar. En particular, el bajo desempeño en
asignaturas fundamentales como Matemáticas y Lengua Portuguesa
representa un desafío crítico.

Este estudio analiza el rendimiento académico en la educación secundaria
mediante un modelo de regresión lineal simple. Para ello, secuenta con
datos de estudiantes, incluyendo calificaciones, características
demográficas, factores sociales y aspectos relacionados con el entorno
escolar, a partir de informes escolares y cuestionarios.

Como aplicación directa de este trabajo, se busca desarrollar una
herramienta interactiva utilizando Shiny y R, que permita modelar y
predecir la calificación final de los estudiantes.

Objetivos:

Ajustar un modelo de regresión lineal simple para predecir la nota final
del estudiante (G3). Este proceso incluirá la exploración de datos
mediante gráficos, la evaluación estadística mediante pruebas t y F,
coeficiente de determinación, el análisis de diagnóstico del modelo, así
como la construcción de intervalos de confianza para los parámetros y
para la estimación de la nota final condicionada a valores de la
variable predictora, además de la generación de predicciones.

Desarrollar una interfaz de usuario con múltiples pestañas que permita
ajustar el modelo, visualizar gráficos y realizar predicciones. Se
implementarán elementos dinámicos, como la selección de variables y la
exclusión de observaciones, para modificar el modelo en tiempo real.
Además, se generarán gráficos interactivos para representar
correlaciones, evaluar el diagnóstico del modelo y visualizar
predicciones con intervalos de confianza.

Elaborar un informe de máximo cuatro páginas que resuma los aspectos más
relevantes del análisis, incluyendo la exploración de datos, el ajuste
del modelo, la inferencia estadística, la evaluación de los supuestos,
la identificación de datos influyentes o atípicos, y la construcción de
intervalos de confianza y predicciones.

Intrucciones:

\begin{enumerate}
\def\labelenumi{\arabic{enumi}.}
\item
  Descargar datos
\item
  Ajustar una regresión lineal simple (RLS) para predecir G3 con la
  variable que mejor la modele. Aplicar estadística descriptivas,
  pruebas de hipótesis, intervalos y todo lo requerido para validar el
  modelo.
\item
  Definir interfaz de usuario (UI): 3.1. Implementar un diseño con
  pestañas: una para la exploración de datos, otra para el resumen del
  modelo, una para diagnóstico y otra para predicción. 3.2. Incluir
  controles dinámicos como para seleccionar la variable predictora y
  para excluir observaciones.
\item
  Definir la lógica del servidor
\end{enumerate}

4.1. Ajustar un modelo de regresión lineal en función de las opciones
seleccionadas en la interfaz.

4.2.Implementar pruebas de supuestos del modelo y presentar los
resultados en tablas interactivas.

4.3. Generar gráficos con para visualizar correlaciones, residuos y
predicciones.

Entregable

Un archivo (.zip) que contenga los datos y un archivo en formato R o
RMarkdown (.Rmd) con el código de la aplicación desarrollada en Shiny y
R.

Un archivo (.zip) que incluya datos, un archivo en formato RMarkdown
(.Rmd) que tenga el informe en HTML creado con RMarkdown de máximo
cuatro páginas (no incluye ni referencias ni anexos) que incluya las
siguientes secciones:

Exploración de datos: Análisis descriptivo mediante gráficos y tablas.

Modelo de regresión lineal: Presentación del modelo y evaluación
estadística a través de pruebas t y F, así como el coeficiente de
determinación.

Análisis de diagnóstico: Evaluación de los supuestos del modelo.

Intervalos de confianza y predicciones: Construcción de intervalos de
confianza para los parámetros del modelo y para la esperanza de la nota
final condicionada a valores de la variable predictora, además de la
generación de predicciones de las notas.

Discusión y conclusiones: Interpretación de los resultados y principales
hallazgos.

Referencias bibliográficas: Fuentes consultadas para el desarrollo del
trabajo.

Anexos: Códigos, gráficos y tablas que respalden el informe.

\end{document}
